The four surfaces tell a reader where a paper sits, but they do not tell a reader whether two
papers' numbers may be placed side by side. For that we need a description of each \emph{side}
of a comparison rather than of each paper. We use five reporting coordinates, each one a
compression of a question the surfaces above already raise.

\begin{itemize}[leftmargin=1.15em,itemsep=1.5pt,topsep=2.5pt,parsep=0pt]
\item[\textbf{C1}] \emph{Evidence state} (\S\ref{sec:evidence}): what the context looks like at
the decision point---supported, distracting, insufficient, conflicting, stale, weakly sourced,
or manipulated---together with the corpus snapshot, retriever, top-$k$, and whether a gold
passage is present at all.
\item[\textbf{C2}] \emph{Origin and capability} (\S\ref{sec:adversarial}): how that state came
about, whether through benign retrieval error, a synthetic stressor, distribution drift, or an
adaptive adversary; for attacks, the adversary's access, persistence, budget unit (documents
injected versus documents retrieved), and whether the score is over all queries or only
triggered ones.
\item[\textbf{C3}] \emph{Intervention} (\S\ref{sec:model}): where the change is made, from the
corpus and retriever through the generator to a verifier module or the evaluation artifact.
\item[\textbf{C4}] \emph{Outcome and unit}: what is scored---answer correctness, retrieval
quality, claim support, attack success, leakage, calibration error, or cost---and per what
unit: a question, a response, or a span.
\item[\textbf{C5}] \emph{Policy and utility} (\S\ref{sec:system}): the recourse available under
uncertainty---answer, verify, re-retrieve, expose the conflict, quarantine, or abstain---and the
costs attached to being wrong, refusing, and following up.
\end{itemize}

\paragraph{A decision procedure for a quoted gap.}
Once both sides are recorded, a gap between them is adjudicated in three steps.
\textbf{(1)}~If C1, C2, C4 and C5 match, the gap \emph{survives} and may be reported as stated.
\textbf{(2)}~If they do not match but one paper reports an operating point matching the other's
configuration, the comparison \emph{narrows} to that point and must be restated there.
\textbf{(3)}~Otherwise the gap is \emph{undefined}: the claim must be confined to the common
subset or withdrawn, not computed across the mismatch. The third outcome is the one current
practice lacks a vocabulary for, and \S\ref{sec:applied} shows how often it is the correct one.

\paragraph{Why C3 is recorded but not gated.}
The intervention coordinate is deliberately excluded from the rule above, and the reason is
definitional rather than empirical: two systems being compared differ in their intervention by
construction, so requiring C3 to match would declare every cross-system comparison invalid.
We nonetheless record C3 and report in \S\ref{sec:tested} exactly what this choice costs.

\paragraph{Relation to existing survey taxonomies.}
Prior surveys organise robust RAG by mechanism---noise types \citep{ragSurvey2024}, trust
dimensions \citep{trustRAGCompass2024}, attack vectors \citep{securityRAGSurvey2026}, conflict
types \citep{xuKnowledgeConflicts2024}, or abstention strategies \citep{wenAbstentionSurvey2025}---and
those structures are effective for finding related work. None of them structures the act of
comparison, and none tests whether paper text carries enough detail for its own categories to be
recovered. Comparability audits of this kind exist outside RAG, in translation, benchmarking and
human evaluation \citep{marie2021credibility,howcroft2020twenty,gehrmann2023cracked}, and our
coordinates adapt long-standing ideas from evidence synthesis, selective prediction and security
evaluation \citep{richardson1995pico,chow1970reject,papernot2018sok} to this setting (Appendix
Table~\ref{tab:main-prior-work-comparison}).
